% Generated draft: requires mathematical and source review.
\documentclass[11pt]{article}
\usepackage[margin=25mm]{geometry}
\usepackage{fontspec}
\setmainfont{Latin Modern Roman}
\usepackage{amsmath,amssymb,mathtools,unicode-math}
\setmathfont{Latin Modern Math}
\usepackage{enumitem,needspace,xurl,hyperref}
\setlength{\parindent}{0pt}
\setlength{\parskip}{5pt}
\setlength{\emergencystretch}{3em}
\newcommand{\N}{\mathbb{N}}
\newcommand{\Z}{\mathbb{Z}}
\newcommand{\Q}{\mathbb{Q}}
\newcommand{\R}{\mathbb{R}}
\newcommand{\C}{\mathbb{C}}
\newcommand{\F}{\mathbb{F}}
\newcommand{\A}{\mathbb{A}}
\newcommand{\PP}{\mathbb{P}}
\newcommand{\E}{\mathbb{E}}
\newcommand{\eps}{\varepsilon}
\newcommand{\dd}{\,\mathrm{d}}
\begin{document}
\section*{UnsolvedMath ID: 700358}
\textbf{Source problem code:} catalogue entry
\subsection{Definitions and mathematical statement}
Let \(\mathfrak L\) be the class of locally compact topological groups with no small subgroups. A topological group \(G\), with identity \(e_G\), has \emph{no small subgroups} if there is an open neighbourhood \(U\) of \(e_G\) such that every subgroup \(K\leq G\) contained in \(U\) is trivial:
\[
K\leq G,\quad K\subseteq U\quad\Longrightarrow\quad K=\{e_G\}.
\]
For a group \(H\), write
\[
Z(H)=\{z\in H:zh=hz\text{ for every }h\in H\}.
\]
A kernel is called discrete when it is discrete in the subspace topology.

\textbf{Requested direct-proof problem.} Study extensions of groups in \(\mathfrak L\) with the objective of giving a direct proof of the following assertion, without using the Gleason--Montgomery--Zippin solution of Hilbert's fifth problem:
\[
\forall G\in\mathfrak L\;\exists H\in\mathfrak L\;\exists\text{ a continuous group homomorphism }\vartheta:G\longrightarrow H
\]
such that
\[
\ker\vartheta\text{ is discrete},\qquad
\vartheta(G)\cap Z(H)=\{e_H\}.
\]
The source specifies no particular type of extension to be studied.

\textbf{Supplied motivation and unresolved qualifications.} The source states that this assertion already follows from the solution of Hilbert's fifth problem, identifying \(\mathfrak L\) with finite-dimensional Lie groups and invoking their local linearity. It further states that a direct proof would substantially shorten the proof of Hilbert's fifth problem, because the adjoint representation of \(H\) would be faithful on \(\vartheta(G)\). These are supplied background claims, not independently verified conclusions of this draft. The record does not specify separation, countability, or connectedness conventions. The intended conventions for the identification with Lie groups, and the hypotheses needed for the stated faithfulness of the adjoint representation, require mathematical review, particularly when \(H\) is disconnected. No such additional hypotheses have been silently imposed.
\subsection{Short English statement}
Dixon asks for a direct proof, through the study of extensions, that every locally compact group with no small subgroups admits a continuous homomorphism into another such group, with discrete kernel and image meeting the target's center only in the identity. The source presents the existence assertion as a known consequence of Hilbert's fifth problem; its supplied open label concerns the requested direct-proof problem. Topological conventions and the adjoint-representation motivation require review.
\subsection{Sources}
\begin{itemize}
\item Supplied local UnsolvedMath registry record: permanent ID \(700358\), exact source code \texttt{catalogue entry}, title \emph{Kourovka Notebook Problem 6.11}. The record labels the problem \texttt{open}; this is supplied registry status, not an independently verified literature assessment.
\item J. D. Dixon, proposer, 1978. \emph{The Kourovka Notebook (Unsolved Problems in Group Theory)}, problem 6.11, sixth issue (1978), as represented in the supplied twenty-first-issue edition (2026), \texttt{arXiv:1401.0300v48}. The supplied collection description names E. I. Khukhro and V. D. Mazurov as editors.
\item D. Montgomery and L. Zippin, \emph{Topological transformation groups}, New York, 1955; cited in the supplied statement for the class of locally compact groups with no small subgroups.
\item V. M. Glushkov, \emph{Uspekhi Matem. Nauk} \textbf{12}, no. 2 (1957), 3--41 (Russian); cited in the supplied statement. No article title is supplied.
\item The supplied statement attributes the solution of Hilbert's fifth problem to Gleason--Montgomery--Zippin, but gives no specific bibliographic citation for that solution. The references and background claims above have not been independently consulted or verified.
\end{itemize}
\end{document}
